\chapter{Arbres et boosting}\label{ch:ml:trees-and-boosting}

Le \omterm{def:ml:trees-and-boosting:boosting}{gradient boosting} par défaut de la société (300 arbres de quinze feuilles) est ajusté sur dix ans d'un panel de 300 actions et obtient $-0.19~\%$ hors échantillon.
Une recherche aléatoire de vingt configurations en retient une qui obtient 0.39~\% sur les
mois de validation et 0.32~\% sur le test~: quatre feuilles, 200 arbres lents. La recherche a fonctionné, pour une fois. Mais en la relançant avec quatre
graines aléatoires différentes, le score de validation des paramètres par défaut varie de 0.00~\% à 0.17~\%, et une recherche qui comparait les
candidats à un tirage unique des paramètres par défaut aurait pu qualifier de succès tout résultat compris entre les deux. Le \omterm{def:ml:trees-and-boosting:boosting}{gradient boosting} est
l'outil de référence du domaine sur les données tabulaires~; ce chapitre le construit à partir d'arbres, mesure son \omterm{def:ml:why-financial-machine-learning-is-different:generalisation}{surapprentissage} sur des données de marché,
et en règle les \omterm{def:ml:why-financial-machine-learning-is-different:sets}{hyperparamètres} en tenant compte du bruit.

\section{Arbres}

\begin{definition}[Arbre de décision]\label{def:ml:trees-and-boosting:tree}
Un \emph{arbre de décision}\index{arbre de décision} partitionne l'espace des variables par une suite de coupures binaires $x_j\le c$,
choisies de manière gloutonne pour réduire au mieux la perte (en régression, la somme des erreurs quadratiques), et prédit dans chaque cellule finale,
ou feuille, la moyenne de la cible de ses lignes d'entraînement. Sa capacité est fixée par sa profondeur et par le nombre minimal de lignes par
feuille.
\end{definition}

\begin{omfigure}
\begin{tikzpicture}[font=\small,level distance=1.25cm,level 1/.style={sibling distance=5.2cm},level 2/.style={sibling distance=2.6cm},
  every node/.style={align=center},split/.style={draw=omDef,fill=omDef!10,rounded corners=2pt,inner sep=3pt},
  leaf/.style={draw=omThm,fill=omThm!10,rounded corners=2pt,inner sep=3pt}]
\node[split] {$c_0 \le 0.1$~?}
  child {node[split] {$c_2 \le -0.4$~?}
    child {node[leaf] {$+0.21~\%$\\ \footnotesize 9\,300 lignes}}
    child {node[leaf] {$-0.18~\%$\\ \footnotesize 17\,100 lignes}}}
  child {node[split] {$c_3 \le 0.5$~?}
    child {node[leaf] {$+0.12~\%$\\ \footnotesize 14\,800 lignes}}
    child {node[leaf] {$+0.47~\%$\\ \footnotesize 6\,800 lignes}}};
\node[font=\footnotesize,anchor=west] at (4.4,-0.1) {oui~: gauche, non~: droite};
\end{tikzpicture}

\omcaption{Arbre de profondeur deux sur des caractéristiques classées, présenté de façon schématique~: chaque coupure envoie une ligne action-mois à gauche ou à droite,
chaque feuille prédit la moyenne des rendements centrés de ses lignes d'entraînement. Les valeurs et effectifs des feuilles sont illustratifs, non
ajustés.}\label{fig:ml:trees:tree}
\end{omfigure}

Un arbre capte les interactions sans qu'on les lui indique (la branche de droite de la \cref{fig:ml:trees:tree} fait intervenir la caractéristique~3
seulement lorsque la caractéristique~0 est élevée) et reste insensible aux transformations monotones des variables. Sa
faiblesse tient à la variance~: les points de coupure et les moyennes des feuilles sont ajustés au bruit. Sur le panel du chapitre (300 actions, 40
caractéristiques classées et risque factoriel, 160 mois d'entraînement, un plafond de 0.77~\% sur les 80 mois de test) un arbre de
profondeur deux obtient 0.24~\% hors échantillon, profondeur cinq $-0.20~\%$, profondeur dix $-3.42~\%$.

\section{Bagging et forêts aléatoires}

\begin{definition}[Bagging, forêt aléatoire, erreur hors sac]\label{def:ml:trees-and-boosting:bagging}
Le \emph{bagging}\index{bagging} (bootstrap aggregation) moyenne les modèles ajustés sur des échantillons bootstrap de l'\omterm{def:ml:why-financial-machine-learning-is-different:sets}{ensemble d'entraînement}.
Une \emph{forêt aléatoire}\index{forêt aléatoire} met en bagging des arbres profonds et, à chaque coupure, ne considère qu'un sous-ensemble aléatoire des
variables, afin de décorréler les arbres. L'\emph{erreur hors sac}\index{erreur hors sac} évalue chaque ligne d'entraînement avec
les arbres dont l'échantillon bootstrap ne la contenait pas.
\end{definition}

\begin{proposition}[Ce qu'apporte la moyenne]\label{prop:ml:trees:bagging}
Si $B$ prédictions ont la même variance $\sigma^2$ et une corrélation deux à deux $\rho$, leur moyenne a pour variance
$\rho\sigma^2 + (1 - \rho)\sigma^2/B$.
\end{proposition}

\begin{proof}
$\Var(\frac1B\sum_bf_b) = \frac1{B^2}(B\sigma^2 + B(B-1)\rho\sigma^2) = \rho\sigma^2 + (1-\rho)\sigma^2/B$.
\end{proof}

Davantage d'arbres n'élimine que le second terme~; le premier disparaît en rendant les arbres moins corrélés, ce à quoi servent les
sous-ensembles aléatoires de variables. Une forêt de 40 arbres (30~\% des variables à chaque coupure, 300 lignes par feuille) obtient 0.20~\%
sur les mois de test. Son R² hors sac vaut 0.42~\%, soit deux fois mieux, et son R² dans l'échantillon 2.76~\%. Les
lignes hors sac ne sont pas hors échantillon sur un panel~: chacune partage son mois, donc ses rendements des facteurs, avec des centaines de lignes en sac
aux caractéristiques voisines. C'est la mise en garde du chapitre~2 sur les scores hors sac, sous une autre forme.

\section{Gradient boosting}

\begin{definition}[Gradient boosting, taux d'apprentissage]\label{def:ml:trees-and-boosting:boosting}
Le \emph{gradient boosting}\index{gradient boosting} construit un modèle additif $F_M(x) = F_0 + \eta\sum_{m=1}^Mh_m(x)$ un
arbre à la fois~: l'arbre $h_m$ est ajusté au gradient négatif de la perte par rapport aux prédictions courantes,
$-\partial\ell(y_i, F_{m-1}(x_i))/\partial F$, et ajouté avec un facteur $\eta\in(0,1]$. Le \emph{taux
d'apprentissage}\index{taux d'apprentissage} $\eta$ est le multiplicateur de pas d'un ajustement itératif~: ici le facteur de rétrécissement de chaque arbre, dans
la descente de gradient, la taille de chaque pas (livre~4, chapitre~24). Suivant la convention du ML, $\eta$ désigne
cette grandeur dans ce livre, comme symbole local.
\end{definition}

\begin{proposition}[Le boosting comme descente de gradient dans l'espace des fonctions]\label{prop:ml:trees:gradient}
Pour la perte quadratique $\ell(y, F) = \frac12(y - F)^2$, le gradient négatif est le résidu $y_i - F_{m-1}(x_i)$, donc chaque
arbre ajuste les résidus courants. Pour toute perte différentiable, un pas $F_m = F_{m-1} + \eta h_m$ avec $h_m$ égal au
gradient négatif aux points d'entraînement diminue la perte empirique pour un $\eta$ suffisamment petit.
\end{proposition}

\begin{proof}
$\partial\frac12(y - F)^2/\partial F = -(y - F)$. Pour la seconde affirmation, $\mathcal L_n(F_{m-1} + \eta h) =
\mathcal L_n(F_{m-1}) - \eta\frac1n\sum_ig_i^2 + O(\eta^2)$ avec $g_i$ le gradient à la ligne $i$, lorsque $h = -g$ aux
points d'entraînement~; un arbre approche $-g$ et le même développement vaut avec $\sum_ig_ih(x_i) < 0$.
\end{proof}

La perte de Huber du livre~4 (chapitre~15) est le choix usuel lorsque les rendements ont des queues épaisses~; les arbres à histogrammes de LightGBM,
construits feuille par feuille, rendent la méthode assez rapide pour ajuster des centaines de modèles par jour (Ke et al.).

\section{Régularisation, arrêt précoce et contraintes de monotonie}

\begin{definition}[Arrêt précoce]\label{def:ml:trees-and-boosting:early}
L'\emph{arrêt précoce}\index{arrêt précoce} (early stopping) arrête un ajustement itératif (plus d'arbres, plus d'époques) à l'itération où la
perte sur un bloc de validation, purgé des lignes d'entraînement, est minimale.
\end{definition}

Le nombre d'arbres et le \omterm{def:ml:trees-and-boosting:boosting}{taux d'apprentissage} se compensent (la \cref{fig:ml:trees:early}). Pour $\eta = 0.1$, le
R² de validation culmine après 34 arbres puis s'effondre~; le modèle arrêté obtient 0.05~\% sur les mois de test,
le modèle poussé à 600 arbres $-3.38~\%$. Pour $\eta = 0.01$, le sommet survient après 233 arbres et est plus haut et plus plat~; le
modèle arrêté obtient 0.26~\% sur les mois de test, et poursuivre jusqu'à 600 arbres coûte peu ($0.03~\%$). Un petit \omterm{def:ml:trees-and-boosting:boosting}{taux d'apprentissage}
est un régularisateur à part entière~: chaque arbre déplace l'ajustement si peu que le bruit doit être ajusté lentement.

\begin{omfigure}
\begin{tikzpicture}
\begin{axis}[width=11cm,height=5.6cm,xlabel={nombre d'arbres},ylabel={$R^2$ de validation (\%)},
  tick label style={font=\small},label style={font=\small},xmin=0,xmax=600,ymin=-1.2,ymax=0.6,grid=major,
  grid style={gray!25},legend columns=2,legend style={font=\footnotesize,at={(0.5,-0.3)},anchor=north,draw=none,
  fill=none,/tikz/every even column/.append style={column sep=8pt}},table/col sep=comma,restrict y to domain=-1.2:1]
\addplot[omThm,very thick] table[x=trees,y=lr01]{figdata/ml/05-trees-and-boosting/early.csv};
\addplot[omDef,very thick] table[x=trees,y=lr001]{figdata/ml/05-trees-and-boosting/early.csv};
\draw[omThm,dashed] (axis cs:34,-1.2) -- (axis cs:34,0.6);
\draw[omDef,dashed] (axis cs:233,-1.2) -- (axis cs:233,0.6);
\legend{taux d'apprentissage 0.1,taux d'apprentissage 0.01}
\end{axis}
\end{tikzpicture}

\omcaption{R-carré de validation (40 mois purgés) des arbres boostés à mesure qu'on les ajoute, à deux \omterm{def:ml:trees-and-boosting:boosting}{taux d'apprentissage}~; en pointillés,
les points d'\omterm{def:ml:trees-and-boosting:early}{arrêt précoce} (34 et 233 arbres). La courbe à $\eta = 0.1$ sort du cadre après environ 340 arbres et atteint
$-2.9\%$ à 600 arbres. Données~: \texttt{ml\_trees.early}.}\label{fig:ml:trees:early}
\end{omfigure}

\begin{definition}[Contrainte de monotonie]\label{def:ml:trees-and-boosting:monotone}
Une \emph{contrainte de monotonie}\index{contrainte de monotonie} impose à un modèle d'être non décroissant (ou non croissant) selon une
variable, les autres étant fixées. Dans un ensemble d'arbres, elle est imposée en n'autorisant une coupure sur cette variable que lorsque la
valeur du fils gauche ne dépasse pas celle du fils droit (ou l'inverse), propagée vers le bas de l'arbre.
\end{definition}

La microstructure connaît souvent un signe que les données sont trop bruitées pour montrer. Dans les sessions de firm.tape, un déséquilibre du flux d'ordres plus fort
ne devrait jamais prévoir un prix milieu plus bas (livre~7, chapitre~8). Des arbres boostés ajustés aux variations du mid à cinq secondes, dans huit
sessions de dix minutes (912 étiquettes sans chevauchement, dix variables de carnet d'ordres), testés sur huit autres, obtiennent un R-carré de
0.333 sans contrainte, 0.364 avec les sept variables de déséquilibre et de flux contraintes à agir positivement. Sans contrainte, une
hausse d'un demi-écart-type du déséquilibre à cinq secondes abaisse la prévision pour 7.0~\% des lignes de test~; avec les
contraintes, aucune (\cref{fig:ml:trees:pdp}). La contrainte est une connaissance a priori utilisée comme régularisation, et elle est aussi
une garantie qu'un comité des risques peut lire. (Les R-carrés sont élevés parce que le mid simulé suit une tendance sur quelques secondes, une
propriété de \texttt{firm.tape} que documente le livre~7~; c'est la comparaison qui compte, non le niveau.)

\begin{omfigure}
\begin{tikzpicture}
\begin{axis}[width=10.5cm,height=5.4cm,xlabel={déséquilibre du flux d'ordres à cinq secondes (unités standardisées)},
  ylabel={prévision moyenne (ticks)},tick label style={font=\small},label style={font=\small},grid=major,
  grid style={gray!25},legend columns=2,legend style={font=\footnotesize,at={(0.5,-0.3)},anchor=north,draw=none,
  fill=none,/tikz/every even column/.append style={column sep=8pt}},table/col sep=comma]
\addplot[omThm,very thick,mark=*,mark size=1.2pt] table[x=x,y=free]{figdata/ml/05-trees-and-boosting/pdp.csv};
\addplot[omDef,very thick,mark=square*,mark size=1.2pt] table[x=x,y=mono]{figdata/ml/05-trees-and-boosting/pdp.csv};
\legend{sans contrainte,monotone}
\end{axis}
\end{tikzpicture}

\omcaption{Dépendance partielle de la prévision à cinq secondes par rapport au déséquilibre du flux d'ordres (la variable balayée de son 2\up{e}
au 98\up{e}~centile, les autres à leurs valeurs observées), pour des arbres boostés sans et avec \omterm{def:ml:trees-and-boosting:monotone}{contraintes de monotonie}. Données~:
\texttt{ml\_trees.partial\_dependence} sur \texttt{firm.tape}.}\label{fig:ml:trees:pdp}
\end{omfigure}

\section{Réglage sous bruit}

Une recherche aléatoire tire vingt configurations (feuilles, \omterm{def:ml:trees-and-boosting:boosting}{taux d'apprentissage}, arbres, lignes par feuille, échantillonnage des lignes et des colonnes, une
pénalité $\ell_2$) et évalue chacune sur le bloc de validation (\cref{fig:ml:trees:tuning}). Leur R-carré de validation
varie de $-0.54~\%$ à 0.39~\%. La configuration gagnante a quatre feuilles, un \omterm{def:ml:trees-and-boosting:boosting}{taux d'apprentissage} de 0.02 et 200 arbres~; réentraînée avec quatre
graines, elle obtient 0.36--0.39~\% (un écart-type de 0.014 points), et 0.32~\% sur les mois de test, soit 0.07 points sous son
score de validation, comme le prédit le chapitre~3 pour une configuration gagnante. Les valeurs par défaut, réentraînées avec les mêmes quatre graines, obtiennent entre
0.00~\% et 0.17~\% (un écart-type de 0.07 points) et $-0.19~\%$ sur les mois de test.

Deux enseignements s'en dégagent. Les configurations complexes sont bruitées autant que mauvaises~: la dispersion d'une graine à l'autre pour les valeurs par défaut est
cinq fois celle de la configuration gagnante, si bien qu'une seule exécution de validation d'un modèle complexe en donne une mauvaise estimation, et une recherche qui
réévalue ses meilleurs candidats avec plusieurs graines est une assurance bon marché. Et la bonne région est un plateau~: six des
vingt configurations se situent à moins de 0.04 point du meilleur. Choisir la configuration la moins complexe, à moins d'un écart-type
entre graines du meilleur (quatre feuilles, 1\,000 lignes par feuille, aucun sous-échantillonnage) donne 0.35~\% en validation et 0.31~\%
sur le test, aussi bon, avec un modèle sept fois moins flexible selon la mesure de la \cref{fig:ml:trees:tuning}.

\begin{omfigure}
\begin{tikzpicture}
\begin{axis}[width=10.5cm,height=5.8cm,xlabel={complexité, $\log_{10}(\text{feuilles}\times\text{arbres}\times\eta/\sqrt{\text{lignes par feuille}})$},
  ylabel={$R^2$ de validation (\%)},tick label style={font=\small},label style={font=\small},grid=major,grid style={gray!25},
  ymin=-0.65,ymax=0.5,xmin=-0.5,xmax=1.8,xtick={-0.4,0,0.4,0.8,1.2,1.6},legend columns=4,legend style={font=\footnotesize,at={(0.5,-0.3)},anchor=north,
  draw=none,fill=none,/tikz/every even column/.append style={column sep=6pt}},table/col sep=comma]
\addplot[only marks,mark=*,mark size=1.8pt,gray] table[x=complexity,y=score]{figdata/ml/05-trees-and-boosting/tuning.csv};
\addplot[only marks,mark=square*,mark size=2.6pt,omThm] table[x=complexity,y=score]{figdata/ml/05-trees-and-boosting/tuning_best.csv};
\addplot[only marks,mark=diamond*,mark size=3pt,omDef] table[x=complexity,y=score]{figdata/ml/05-trees-and-boosting/tuning_plateau.csv};
\addplot[only marks,mark=triangle*,mark size=2.4pt,omProp] table[x=complexity,y=score]{figdata/ml/05-trees-and-boosting/tuning_default.csv};
\legend{autres configurations,meilleure,choix du plateau,valeurs par défaut (4 graines)}
\end{axis}
\end{tikzpicture}

\omcaption{Vingt configurations aléatoires d'arbres boostés sur les mois de validation, rapportées à une mesure grossière de leur
flexibilité~; les valeurs par défaut réentraînées avec quatre graines. Données~: \texttt{ml\_trees.tuning}.}\label{fig:ml:trees:tuning}
\end{omfigure}

\begin{method}[Réglage des arbres boostés sur données de marché]\label{met:ml:trees:tuning}
\begin{enumerate}
\item Fixer un bloc de validation purgé (ou des plis purgés) et ne jamais régler sur le test.
\item Partir d'un petit \omterm{def:ml:trees-and-boosting:boosting}{taux d'apprentissage} (0.01--0.03) avec \omterm{def:ml:trees-and-boosting:early}{arrêt précoce} (early stopping), peu de feuilles et beaucoup de lignes par feuille.
\item Chercher au hasard parmi les termes de régularisation~; réévaluer les quelques meilleurs candidats avec plusieurs graines.
\item Choisir sur le plateau (la moins flexible à moins d'un écart-type entre graines du meilleur), et non au sommet.
\item Déclarer des \omterm{def:ml:trees-and-boosting:monotone}{contraintes de monotonie} partout où l'économie fixe un signe.
\end{enumerate}
\end{method}

\section{Tutoriel~: réglage dans le brouillard}\label{tut:ml:trees-and-boosting:tut}

\begin{tutorial}
\textbf{Objectif.} Ajuster des arbres, une forêt et des arbres boostés aux données de panel~; arrêter le boosting tôt à deux \omterm{def:ml:trees-and-boosting:boosting}{taux d'apprentissage}~; lancer une
recherche aléatoire et choisir sur son plateau~; contraindre les variables du carnet d'ordres sur \texttt{firm.tape}. \textbf{État final~:}
les chiffres du chapitre et les \cref{fig:ml:trees:early,fig:ml:trees:pdp,fig:ml:trees:tuning}.
\begin{tutsteps}
\item \textbf{\omterm{def:ml:trees-and-boosting:early}{Arrêt précoce}} par les callbacks de LightGBM, en conservant la courbe de validation.
\omcode{code/firm/gbdt/firm_gbdt.py}{38}{48}{Arrêt précoce sur un bloc de validation purgé.}\label{lst:ml:trees:early}
\item \textbf{La règle du plateau et l'export du modèle}~: les arbres ajustés deviennent un dictionnaire qu'une fonction NumPy
évalue exactement (le chapitre~26 compile le même dictionnaire).
\omcode{code/firm/gbdt/firm_gbdt.py}{63}{98}{Choisir sur le plateau~; évaluer les arbres exportés.}\label{lst:ml:trees:dump}
\item \textbf{Exécuter} \texttt{ml\_trees.family()}, \texttt{early()}, \texttt{tuning()}, \texttt{monotone()} et
\texttt{fig\_trees.py}.
\end{tutsteps}
\textbf{Pour aller plus loin.} Contraindre le déséquilibre du flux d'ordres avec le mauvais signe et mesurer le coût~; réévaluer chaque
configuration de la recherche avec quatre graines et voir si la configuration gagnante change.
\end{tutorial}

\section{Construction~: les arbres boostés de la firme}\label{bld:ml:trees-and-boosting:gbdt}

\begin{build}
\bfield{Rôle} Une façon d'ajuster, d'arrêter, de contraindre, d'assembler, de choisir et d'exporter des arbres boostés par gradient.
\bfield{Interface} \texttt{make(params, monotone, names, seed)}, \texttt{DEFAULTS},
\texttt{fit\_early\_stop(params, X, y, Xv, yv, max\_trees, patience, monotone, names, seed)}, \texttt{SeedEnsemble(params, seeds, monotone, names)},
\texttt{plateau\_select(scores, se, complexity)}, \texttt{to\_dict(model)}, \texttt{predict\_dict(d, X)}.
\bfield{Règles} Ajustements déterministes, monothread~; le bloc de validation est purgé par l'appelant~; les contraintes sont
déclarées par nom de variable, jamais par position de colonne~; les dictionnaires exportés reproduisent exactement les prédictions, y compris
les valeurs manquantes.
\bfield{Tests d'acceptation} \texttt{code/firm/gbdt/tests/}~: deux ajustements donnent des prédictions identiques~; l'évaluateur de dictionnaires
concorde avec LightGBM sur les valeurs manquantes~; l'\omterm{def:ml:trees-and-boosting:early}{arrêt précoce} conserve la meilleure itération et l'export ne contient que celle-ci~; une variable monotone
ne fait jamais baisser la prédiction~; les ensembles de graines moyennent~; la règle du plateau retient le candidat le plus simple proche du meilleur.
\bfield{Pour aller plus loin} Pertes en quantiles et perte de Huber à travers la même interface~; un réajustement avec initialisation à chaud qui ajoute des arbres pour de nouveaux
mois (chapitre~12).
\end{build}

\begin{omsources}
\item L.~Breiman, ``Random forests'', \emph{Machine Learning} 45, 2001.
\item J.~H.~Friedman, ``Greedy function approximation: a gradient boosting machine'', \emph{Annals of Statistics} 29(5),
2001.
\item G.~Ke et al., ``LightGBM: a highly efficient gradient boosting decision tree'', \emph{Advances in Neural Information
Processing Systems} 30, 2017.
\item T.~Chen et C.~Guestrin, ``XGBoost: a scalable tree boosting system'', \emph{KDD}, 2016.
\item L.~Grinsztajn, E.~Oyallon et G.~Varoquaux, ``Why do tree-based models still outperform deep learning on typical
tabular data?'', \emph{NeurIPS Datasets and Benchmarks}, 2022.
\end{omsources}

\section{Exercices}

\begin{exercise}[$\star$]\label{exo:ml:trees-and-boosting:1}
Cinquante arbres ont chacun une variance de prédiction $\sigma^2$ et une corrélation deux à deux de 0.3. Calculer la variance de leur
moyenne, puis celle obtenue avec une infinité d'arbres.
\end{exercise}

\begin{exercise}[$\star$]\label{exo:ml:trees-and-boosting:2}
Un arbre partage 1\,000 lignes en deux feuilles de 400 et 600 lignes dont les cibles moyennes valent $-0.1$ et $+0.2$. La
moyenne globale vaut $0.08$. Calculer de combien le partage réduit la somme des carrés des erreurs.
\end{exercise}

\begin{exercise}[$\star$]\label{exo:ml:trees-and-boosting:3}
Le boosting avec perte quadratique part de $F_0 = 0$~; une ligne a pour cible 1.0 et le premier arbre prédit 0.8 pour elle. Avec
$\eta = 0.1$, calculer le résidu de cette ligne après la première étape, puis celui obtenu après dix étapes si chaque arbre, le premier
compris, prédisait exactement son résidu courant.
\end{exercise}

\begin{exercise}[$\star\star$]\label{exo:ml:trees-and-boosting:4}
Expliquer pourquoi le R-carré hors sac de la forêt (0.42~\%) vaut deux fois son R-carré de test (0.20~\%) sur le panel.
\end{exercise}

\begin{exercise}[$\star\star$]\label{exo:ml:trees-and-boosting:5}
Le score de validation des paramètres par défaut sur quatre graines vaut 0.17, 0.11, 0.09 et 0.00~\%. Un candidat obtient 0.14~\% sur une seule graine.
Dire s'il est meilleur que les paramètres par défaut, et que faire avant de trancher~?
\end{exercise}

\begin{exercise}[$\star\star$]\label{exo:ml:trees-and-boosting:6}
\emph{Trouver la faille.} «~On a utilisé l'\omterm{def:ml:trees-and-boosting:early}{arrêt précoce} sur les mois de test pour choisir le nombre d'arbres, et on rapporte le
R-carré de test obtenu à ce point.~»
\end{exercise}

\begin{exercise}[$\star\star\star$]\label{exo:ml:trees-and-boosting:7}
\emph{Programmation.} Dans \texttt{ml\_trees.monotone}, imposer que le déséquilibre du flux d'ordres à cinq secondes agisse \emph{négativement}
(toutes les autres contraintes comme dans le chapitre). Rapporter le R-carré de test et expliquer le résultat.
\end{exercise}

\begin{exercise}[$\star\star\star$]\label{exo:ml:trees-and-boosting:8}
Montrer qu'avec une perte quadratique et un \omterm{def:ml:trees-and-boosting:boosting}{taux d'apprentissage} $\eta$, si chaque arbre ajustait exactement les résidus, le résidu
d'une ligne après $M$ arbres vaudrait $(1 - \eta)^M$ fois sa cible~; expliquer pourquoi cela rend l'\omterm{def:ml:trees-and-boosting:early}{arrêt précoce} et un petit \omterm{def:ml:trees-and-boosting:boosting}{taux
d'apprentissage} interchangeables.
\end{exercise}

\section{Problème~: Réglage dans le brouillard}

\begin{problem}[{Problème du week-end --- vingt configurations et un plateau}]\label{pb:ml:trees-and-boosting:1}
Le panel du chapitre~: 300 actions, quarante caractéristiques à risque factoriel, 160 mois d'entraînement (les 40 derniers servant de
bloc de validation purgé lors du réglage), 80 mois de test.

\medskip\noindent\textbf{Partie I --- Arbres isolés et forêts.}
\begin{enumerate}
\item Donner le score hors échantillon des arbres de profondeur 2, 5 et 10, et expliquer pourquoi cet ordre.
\item Donner le score de la forêt hors échantillon, hors sac et dans l'échantillon.
\item Expliquer pourquoi le score hors sac est optimiste sur un panel.
\item Dire ce que dit la \cref{prop:ml:trees:bagging} sur ce que davantage d'arbres peuvent et ne peuvent pas faire.
\end{enumerate}

\medskip\noindent\textbf{Partie II --- Boosting et \omterm{def:ml:trees-and-boosting:early}{arrêt précoce}.}
\begin{enumerate}[resume]
\item Donner le score des paramètres par défaut sur les mois de test.
\item Dire où s'arrête l'\omterm{def:ml:trees-and-boosting:early}{arrêt précoce} pour $\eta = 0.1$ et $0.01$, et donner le score des modèles ainsi arrêtés.
\item Évaluer le coût qu'aurait eu un entraînement jusqu'à 600 arbres pour chaque \omterm{def:ml:trees-and-boosting:boosting}{taux d'apprentissage}.
\item Expliquer pourquoi un petit \omterm{def:ml:trees-and-boosting:boosting}{taux d'apprentissage} est une forme de régularisation.
\end{enumerate}

\medskip\noindent\textbf{Partie III --- La recherche.}
\begin{enumerate}[resume]
\item Donner l'étendue des scores de validation couverts par les vingt configurations.
\item Identifier la configuration gagnante, et donner son score avec quatre graines et sur les mois de test.
\item Comparer la dispersion des paramètres par défaut selon les graines à celle de la configuration gagnante.
\item Dire quelle configuration retient la règle du plateau, et donner son score.
\end{enumerate}

\medskip\noindent\textbf{Partie IV --- Contraintes et verdict.}
\begin{enumerate}[resume]
\item Décrire l'effet des \omterm{def:ml:trees-and-boosting:monotone}{contraintes de monotonie} sur le R-carré du modèle de carnet d'ordres et sur ses violations.
\item Expliquer pourquoi une contrainte peut améliorer un score hors échantillon.
\item Énoncer le \emph{résultat nommé}~: l'écart-type selon les graines du score de validation des paramètres par défaut et de celui du gagnant,
comparé au gain du réglage, ainsi que l'écart entre validation et test du gagnant.
\item Dire quelle configuration mettre en production, et pourquoi.
\item Dire combien de graines utiliser pour réévaluer les meilleurs candidats d'une recherche.
\item Indiquer ce que le journal de recherche doit consigner sur la recherche.
\item Dire ce qui changerait avec dix fois plus d'actions.
\item En une phrase, dire comment régler les arbres boostés sur des données de marché.
\end{enumerate}
\end{problem}

\section{Questions d'entretien}

\begin{interviewq}[$\star$ \iqroles{chercheur, ingénieur ML}]\label{iq:ml:trees-and-boosting:1}
Expliquez la différence entre \omterm{def:ml:trees-and-boosting:bagging}{bagging} et boosting.
\end{interviewq}

\begin{interviewq}[$\star$ \iqroles{ingénieur ML}]\label{iq:ml:trees-and-boosting:2}
Quels \omterm{def:ml:why-financial-machine-learning-is-different:sets}{hyperparamètres} d'un modèle de \omterm{def:ml:trees-and-boosting:boosting}{gradient boosting} régleriez-vous en premier sur des données bruitées, et dans quel sens~?
\end{interviewq}

\begin{interviewq}[$\star\star$ \iqroles{chercheur}]\label{iq:ml:trees-and-boosting:3}
Pourquoi l'\omterm{def:ml:trees-and-boosting:bagging}{erreur hors sac} d'une \omterm{def:ml:trees-and-boosting:bagging}{forêt aléatoire} peut-elle être trompeuse sur des données de panel financières~?
\end{interviewq}

\begin{interviewq}[$\star\star$ \iqroles{chercheur, trader}]\label{iq:ml:trees-and-boosting:4}
Qu'est-ce qu'une \omterm{def:ml:trees-and-boosting:monotone}{contrainte de monotonie}~? Citez deux variables d'un modèle de trading que vous contraindriez.
\end{interviewq}

\begin{interviewq}[$\star\star$ \iqroles{ingénieur ML}]\label{iq:ml:trees-and-boosting:5}
Deux configurations de boosting diffèrent de 0.03 point de R-carré en validation. Comment trancheriez-vous entre elles~?
\end{interviewq}

\begin{interviewq}[$\star\star\star$ \iqroles{chercheur, ingénieur ML}]\label{iq:ml:trees-and-boosting:6}
Démontrez pourquoi le boosting avec perte quadratique ajuste les résidus, et ce qui change avec la perte de Huber.
\end{interviewq}
